\section{Introduction}
Algorithmic paradigms provide structured methodologies to formulate and solve complex computational challenges with rigorous guarantees on optimality and runtime complexity. Among these paradigms, \textit{greedy algorithms} construct solutions iteratively by making locally optimal decisions without backtracking, while \textit{divide-and-conquer} techniques recursively decompose problems into independent subproblems, solve them, and merge their results into a global solution.

In this work, we present an in-depth investigation of these paradigms applied to two critical domains:
\begin{enumerate}
    \item \textbf{Optimal Satellite and Cloud Resource Scheduling (Greedy):} Maximizing observation mission throughput under temporal non-overlap constraints, modeled as an optimal interval scheduling problem.
    \item \textbf{Air Traffic Proximity Warning and Drone Collision Detection (Divide-and-Conquer):} Determining minimum Euclidean spatial separations within dense airspace, modeled as the 2D closest pair problem.
\end{enumerate}

For both problems, we establish formal mathematical abstractions, specify optimal algorithms with asymptotic complexity analyses, prove correctness via formal inductive and exchange arguments, explain the operational mechanisms in domain-specific terminology, and provide reproducible empirical benchmarks validating theoretical predictions.

